\documentclass[11pt]{article}

% ---------- Packages ----------
\usepackage[a4paper,margin=1in]{geometry}
\usepackage{amsmath,amssymb,amsthm,mathtools}
\usepackage{hyperref}
\usepackage{enumitem}
\usepackage{booktabs}
\usepackage{array}
\usepackage{caption}

% ---------- Theorem environments ----------
\newtheorem{theorem}{Theorem}[section]
\newtheorem{proposition}[theorem]{Proposition}
\newtheorem{lemma}[theorem]{Lemma}
\newtheorem{corollary}[theorem]{Corollary}

\theoremstyle{definition}
\newtheorem{definition}[theorem]{Definition}
\newtheorem{example}[theorem]{Example}

\theoremstyle{remark}
\newtheorem{remark}[theorem]{Remark}
\newtheorem*{problem}{Open Problem}

% ---------- Commands ----------
\newcommand{\St}{\mathrm{St}}
\newcommand{\Clop}{\mathrm{Clop}}
\newcommand{\CO}{\mathrm{CO}}
\newcommand{\tr}{\mathrm{tr}}
\newcommand{\PRlat}{\mathrm{PR}_{\mathrm{lat}}}
\newcommand{\PRdual}{\mathrm{PR}_{\mathrm{dual}}}

% ---------- Title ----------
\title{$\sigma$-Additive Probability on Orthomodular Lattices:\\
A Survey and an Open Problem}
\author{Patrick S.\ Mahon}
\date{Survey note --- June 2026}

\begin{document}
\maketitle

\begin{abstract}
The passage from a finitely coherent assignment of probabilities to a
countably additive measure is, in the Boolean case, governed by a single
condition (countable additivity) and executed by classical machinery
(Carath\'{e}odory; Loomis--Sikorski; Stone duality). On a
\emph{non-distributive} orthomodular lattice (OML) --- the algebra of
propositions of a quantum system --- this machinery breaks, and the
question of when finite coherence forces countable behaviour is open. This
note surveys the problem at the level needed to begin work on it. We fix
the apparatus (orthomodular lattices, the gap between orthogonality and
meet-zero, the state/valuation/charge ladder, the McDonald--Bimb\'{o} dual
space), separate the question into two independent axes --- \emph{extension}
and \emph{descent} --- review what is known and what is ruled out, and
isolate the one residue that remains genuinely open: whether a concrete,
non-Boolean, infinite $\sigma$-complete OML admits a Loomis--Sikorski-type
representation. We close with the first concrete move. This note is a guide
to the companion research note \cite{mahonOMLnote}; the underlying
verification scripts are referenced where claims are checked.
\end{abstract}

% ==================================================================
\section{Introduction}
\label{sec:intro}
% ==================================================================

All measurement is finite, yet the frameworks that organise empirical
knowledge --- probability theory, quantum mechanics --- rest on infinite
structures. The passage from finite observation to infinite structure is
not automatic; it requires a commitment no finite body of evidence can
compel. In probability theory the locus of that commitment is
\emph{countable additivity}: that mass cannot persist along a vanishing
sequence of events. De~Finetti~\cite{deFinetti1974} held that only finite
additivity is empirically grounded; Kolmogorov~\cite{Kolmogorov} adopted
$\sigma$-additivity as an axiom.

In quantum theory the same tension recurs in a different guise. The algebra
of propositions about a quantum system is not Boolean: incompatible
observables cannot be jointly determined, and the closed subspaces of a
Hilbert space $H$ form an \emph{orthomodular lattice} $L(H)$ in which
distributivity may fail. The classical extension machinery ---
Carath\'{e}odory's outer measure, which needs a Boolean algebra of sets ---
does not apply. For the single lattice $L(H)$ this is rescued by Gleason's
rigidity theorem~\cite{Gleason1957}, but Hilbert space is special: for
general OMLs no analogue is known.

This note asks the shared question --- \emph{when does finite coherence
force countable behaviour?} --- in the OML setting, and surveys it to the
point where one can begin work. The contribution of the companion
note~\cite{mahonOMLnote} is largely a clarification: the question, vague as
usually posed, splits into two independent axes, one of which turns out to
be classical and settled, the other genuinely open. The survey here is
organised around that split.

\paragraph{The relational standpoint.} We work in the spirit of the
``structure from observation'' programme~\cite{mahon2025}: rather than
begin with a space $\Omega$ of outcomes and ask which $\sigma$-algebras it
supports, begin with the propositions and their entailment order, and ask
what that relational structure alone determines. For non-Boolean algebras
the natural dual is not a space of points but the McDonald--Bimb\'{o}
space of \emph{filters} (\S\ref{sec:dual}). This is what makes the open
problem worth its difficulty: a positive answer would be a
\emph{point-free} (relational) $\sigma$-additive probability theory on a
non-distributive lattice --- the non-Boolean analogue of localic measure
theory (\S\ref{sec:problem}).

% ==================================================================
\section{Preliminaries}
\label{sec:prelim}
% ==================================================================

We collect the lattice-theoretic apparatus. All notions are standard
except \emph{valuation}, where we adopt a local convention flagged below.
The elementary laws are exactly those verified for $\mathrm{MO}_3$ in the
companion scripts (\texttt{verification/mo3\_extension.py}).

\subsection{Orthomodular lattices, and the orthogonal/meet-zero gap}
\label{sec:oml}

\begin{definition}[Orthocomplemented lattice]\label{def:ocl}
A \emph{bounded lattice} is a partially ordered set in which every pair
$a,b$ has a join $a \vee b$ and meet $a \wedge b$, with greatest element
$\mathbf 1$ and least element $\mathbf 0$. An \emph{orthocomplementation}
is a map $a \mapsto a^\perp$ satisfying, for all $a,b$:
\begin{enumerate}[label=(\roman*), noitemsep]
  \item $a \wedge a^\perp = \mathbf 0$ and $a \vee a^\perp = \mathbf 1$
    (complement),
  \item $a^{\perp\perp} = a$ (involution),
  \item $a \leq b \Rightarrow b^\perp \leq a^\perp$ (order-reversing).
\end{enumerate}
\end{definition}

\begin{definition}[Orthogonal; meet-zero]\label{def:ortho}
Elements $a,b$ are \emph{orthogonal}, written $a \perp b$, if
$a \leq b^\perp$ (equivalently $b \leq a^\perp$). They are
\emph{meet-zero} if $a \wedge b = \mathbf 0$. Orthogonality always implies
meet-zero; the converse holds in every Boolean algebra but \textbf{fails}
once the lattice is non-distributive. This gap drives everything below.
\end{definition}

\begin{definition}[Orthomodular lattice]\label{def:oml}
An orthocomplemented lattice $A$ is \emph{orthomodular} if
\begin{equation}\label{eq:oml-law}
  a \leq b \quad\Longrightarrow\quad b = a \vee (a^\perp \wedge b).
\end{equation}
A \emph{Boolean algebra} is the special case where $\wedge$ distributes
over $\vee$; every Boolean algebra is orthomodular, not conversely. The
motivating non-Boolean example is $L(H)$, the closed subspaces of a
Hilbert space with $a^\perp$ the orthogonal complement.
\end{definition}

\begin{example}[$\mathrm{MO}_n$, and the gap made concrete]\label{ex:mo}
$\mathrm{MO}_n$ is the lattice with $\mathbf 0$, $\mathbf 1$, and $n$
complementary pairs of incomparable atoms, all sharing only $\mathbf 0$
and $\mathbf 1$. It is orthomodular and non-distributive ($n \geq 2$): any
two distinct atoms $a,b$ have $a \wedge b = \mathbf 0$ yet are
\emph{not} orthogonal, so meet-zero $\neq$ orthogonal already here.
$\mathrm{MO}_3$ is the smallest non-distributive OML.
\end{example}

\begin{definition}[OMP; concrete logic]\label{def:omp}
An \emph{orthomodular poset} (OMP) weakens Definition~\ref{def:oml}: joins
$a \vee b$ are required only for \emph{orthogonal} pairs. An OMP is a
\emph{concrete logic} (equivalently \emph{set-representable}) if it embeds
into $(\mathcal P(X), \subseteq, X \setminus(-))$ for some set $X$,
preserving existing orthogonal joins.
\end{definition}

\begin{remark}[Concreteness does not close the gap]\label{rem:concrete}
A set-representable OMP needs only closure under complement and
\emph{disjoint} union~\cite{BuresovaPtak2024}; intersection-closure would
force a field of sets, hence Boolean. The lattice meet $a \wedge b$ is the
greatest $L$-member inside $A \cap B$, so orthogonality
($A \cap B = \varnothing$) implies meet-zero, but not conversely. Indeed
$\mathrm{MO}_3$ is \emph{itself} concrete (it carries an order-determining
set of two-valued states; Gudder), so the flagship meet-zero-$\neq$-
orthogonal example is a concrete logic. The condition that closes the gap
is not concreteness but a richness/atomicity condition (meet-zero
coincides with orthogonality iff every nonempty $A \cap B$ contains a
nonzero $L$-member); the abstract form $a \wedge b = 0 \Rightarrow a \perp
b$ is Tkadlec's \emph{Boolean orthoposet} condition (which does
\emph{not} mean Boolean \emph{algebra}). Verified:
\texttt{verification/concrete\_meetzero\_vs\_orthogonal.py}.
\end{remark}

\subsection{The state/valuation/charge ladder}
\label{sec:ladder}

\begin{definition}[State; valuation; charge-extendible]\label{def:ladder}
Let $A$ be an OML. A \emph{state} is a map $s : A \to [0,1]$ with
$s(\mathbf 1)=1$ that is additive on \emph{orthogonal} pairs:
$s(a \vee b) = s(a) + s(b)$ whenever $a \perp b$. Three strengthenings,
strictly increasing once $A$ is non-distributive:
\begin{enumerate}[label=(\arabic*), noitemsep]
  \item \emph{State} --- additive on orthogonal pairs.
  \item \emph{Valuation} (local convention) --- additive on every
    \emph{meet-zero} pair $a \wedge b = \mathbf 0$ (binary, strictly
    stronger). \textbf{Warning:} not the standard lattice-theoretic
    valuation $v(a)+v(b)=v(a\wedge b)+v(a \vee b)$.
  \item \emph{Charge-extendible} --- the $n$-ary condition
    $\sum_i s(a_i) \leq 1$ for every pairwise-meet-zero family $\{a_i\}$
    (strictly stronger again). This is the \emph{extension} condition of
    \S\ref{sec:axes}.
\end{enumerate}
\end{definition}

\begin{example}[The ladder is strict]\label{ex:strict}
On $\mathrm{MO}_3$ with complementary pairs $(p,p^\perp),(q,q^\perp),
(r,r^\perp)$, the maximally mixed $s \equiv \tfrac12$ is a valuation
(clears~(2)) yet fails~(3) at $\{p,q,r\}$:
$\tfrac12+\tfrac12+\tfrac12 = \tfrac32 > 1$. So no state on
$\mathrm{MO}_3$ is charge-extendible. \emph{(\texttt{mo3\_extension.py}.)}
\end{example}

\subsection{The McDonald--Bimb\'{o} dual space}
\label{sec:dual}

\begin{definition}[Dual space; representation map]\label{def:dual}
The McDonald--Bimb\'{o} duality~\cite{McDonaldBimbo2023} assigns to an OML
$A$ a compact space
\[
  S_0(A) = \bigl(F(A), \subseteq, \perp_A, P(A), T(S)\bigr),
\]
where $F(A)$ is the set of \emph{filters} of $A$ and $P(A) \subseteq F(A)$
the \emph{principal} filters --- the ``physical points,'' the realisation
datum (canonical, unlike the non-canonical pure points of the Boolean
case). For $a \in A$ set $h(a) = \{x \in F(A) : a \in x\}$; then $h$ is an
OML isomorphism onto the $\perp$-stable clopens $\CO(S_0(A))^\dagger$.
Define $\mu(h(a)) = s(a)$.
\end{definition}

\begin{remark}[The duality is finitary --- the wall]\label{rem:finitary}
The identity $h(\bigvee_n a_n) = (\bigcup_n h(a_n))^{\perp\perp}$ holds for
\emph{finite} joins but fails for countable ones: finitary OML
homomorphisms need not preserve infinite suprema. This finitariness is the
wall the open problem of \S\ref{sec:problem} runs into. The rival OML
duality of Cannon--D\"{o}ring~\cite{Cannon2013} is also finitary and
carries no measure; we use McDonald--Bimb\'{o} because it carries $P(A)$ as
explicit structure. Verified against the primary source:
\texttt{verification/mb\_primeness\_check.md}.
\end{remark}

\subsection{The Boolean engine: Loomis--Sikorski}
\label{sec:ls}

In the Boolean case, the descent argument
(``$\sigma$-additive $\Leftrightarrow$ the measure concentrates on the
physical points'') runs on a \emph{Loomis--Sikorski} representation: a
$\sigma$-complete Boolean algebra is a $\sigma$-tribe of sets modulo a
$\sigma$-ideal. There is \emph{no} OML analogue, and supplying one --- or
proving none exists --- is the open problem of \S\ref{sec:problem}. The
slogan: the Boolean proof has an engine; the OML proof needs one and does
not have it.

% ==================================================================
\section{Two axes: extension and descent}
\label{sec:axes}
% ==================================================================

The headline question --- \emph{when does a state $s$ on an OML $A$ extend
to a $\sigma$-additive measure on $S_0(A)$?} --- splits into two
\textbf{independent} axes that the Boolean case fuses. A state $s$ on $A$
faces:

\begin{center}
\renewcommand{\arraystretch}{1.35}
\begin{tabular}{@{}p{2.3cm} p{5.3cm} p{5.3cm}@{}}
\toprule
 & \textbf{Extension axis} & \textbf{Descent axis} \\
\midrule
\textbf{Question}
  & Does $\mu$ extend to a finitely additive charge on the \emph{full}
    Boolean clopen algebra of $S_0(A)$?
  & Once extended, does the measure concentrate on the physical points
    $P(A)$? \\
\textbf{Governed by}
  & the meet-zero/orthogonal gap (\S\ref{sec:oml})
  & $\sigma$-additivity --- needs the missing engine (\S\ref{sec:ls}) \\
\textbf{Status}
  & \textbf{settled}: Horn--Tarski / Pitowsky feasibility; finite case
    $=$ decidable LP; can fail for every state.
  & \textbf{open}. \emph{This is the residue.} \\
\bottomrule
\end{tabular}
\end{center}

\noindent Two observations pin down the extension axis.

\begin{proposition}[Finite case]\label{prop:finite}
For finite $A$ the extension condition is a decidable linear program
(Horn--Tarski~\cite{HornTarski1948} / Pitowsky correlation-polytope
feasibility~\cite{Pitowsky1989}); since $h$ preserves meets, meet-zero
elements map to disjoint clopens and a charge must satisfy
$\sum_i s(a_i) \leq 1$ over pairwise-meet-zero families. It can fail for
\emph{every} state ($\mathrm{MO}_3$, Example~\ref{ex:strict}).
\end{proposition}

\begin{proposition}[Infinite $L(H)$, normal states]\label{prop:lh}
Let $\dim H = \infty$. Any state $s$ on $L(H)$ with $s(e) > 0$ for some
finite-dimensional $e$ fails to extend to a charge on $S_0(L(H))$; in
particular no normal (Gleason) state extends.
\end{proposition}

\begin{proof}[Proof sketch]
Fix a $2$-plane $e$. For $k$ distinct lines $p_1,\dots,p_k \subset e$ the
$2k$ lines $\{p_i, p_i^\perp\}$ are pairwise meet-zero, so $h$ sends them
to disjoint clopens and orthoadditivity forces
$\sum_i [s(p_i)+s(p_i^\perp)] = k\,s(e) \leq 1$ for every $k$. If
$s(e)>0$, take $k > 1/s(e)$. The argument is finitary;
$\sigma$-completeness of the lattice is irrelevant.
(\texttt{verification/lh\_infinite\_extension.py};
\texttt{lh\_singular\_dichotomy.md}.)
\end{proof}

By the Takesaki normal/singular decomposition (lifted via Mackey--Gleason
/ Bunce--Wright~\cite{BunceWright1992}, valid as $B(H)$ has no type
$\mathrm{I}_2$ summand), the only states escaping
Proposition~\ref{prop:lh} are the \emph{singular} ones ($s(e)=0$ for all
finite-dim $e$; e.g.\ ultrafilter vector states, or pullbacks of
Calkin-algebra states). These form a clean dichotomy with the normal
states. So on the extension axis the sole survivors are the singular
states of $L(H)$ --- the sliver of \S\ref{sec:first}. The genuine open
problem lives on the descent axis (\S\ref{sec:problem}).

% ==================================================================
\section{What is known}
\label{sec:known}
% ==================================================================

\subsection{The Boolean benchmark}
In the Boolean case the Kelley--Vladimirov--Pt\'{a}k characterisation
(Fremlin~\cite{Fremlin3}, Thm 391D) is complete: a Boolean algebra carries
a strictly positive $\sigma$-additive measure iff it is Dedekind
$\sigma$-complete, weakly $(\sigma,\infty)$-distributive, and chargeable.
Each condition has \emph{no} clean OML analogue: weak
$(\sigma,\infty)$-distributivity relies on distributivity of $\bigwedge$
over $\bigvee$; chargeability has no structural characterisation on OMLs.

\subsection{Positive and obstructive OML results}
The positive results require structure rich enough to approximate
Hilbert-space geometry: Gleason~\cite{Gleason1957} ($L(H)$,
$\dim \geq 3$; $\sigma$-additivity \emph{on the lattice}),
Bunce--Wright~\cite{BunceWright1992} (JBW projection lattices),
Chetcuti--Dvure\v{c}enskij~\cite{ChetcutiDvurecenskij2005}. The
obstructive results show the Boolean machinery cannot be transplanted:
Pt\'{a}k--Pulmannov\'{a}~\cite{PtakPulmannova1991} (conditions forcing
$\sigma$-additivity collapse the lattice to Boolean),
Navara~\cite{Navara1992} (regularity does not force $\sigma$-additivity),
Pt\'{a}k~\cite{Ptak1987} (exotic state spaces via Greechie pasting).

\subsection{Three blocked routes to a $\sigma$-OML engine}
\label{sec:routes}
The descent argument needs an OML Loomis--Sikorski. The three known routes
are all blocked.
\begin{enumerate}[label=(\roman*)]
  \item \textbf{RDP / effect-algebra.} Loomis--Sikorski holds for
    $\sigma$-MV algebras~\cite{Dvurecenskij2000} and monotone
    $\sigma$-effect algebras \emph{with} the Riesz Decomposition Property
    \cite{Dvurecenskij2005}; but a lattice effect algebra has RDP iff it
    is MV, and OML $\cap$ MV $=$ Boolean (Kalmbach~\cite{Kalmbach1983}).
    So $\text{OML}+\text{RDP} \Leftrightarrow \text{Boolean}$.
  \item \textbf{MacNeille completion.} The MacNeille completion of an OML
    need not be orthomodular (Harding~\cite{Harding1991}); one cannot
    complete to absorb countable joins.
  \item \textbf{$\sigma$-Stone duality.} None exists. McDonald--Bimb\'{o}
    is finitary (Remark~\ref{rem:finitary}); Freytes'
    theory~\cite{Freytes2020} of $\sigma$-complete OMLs is equational, with
    no set/tribe representation.
\end{enumerate}

\subsection{Prior art on point-free quantum probability}
\label{sec:priorart}
Every existing point-free or directed-context construction lands on one
side of the $\sigma$-additivity-on-a-non-distributive-structure line.
\begin{description}[style=nextline, font=\normalfont\bfseries,
    labelwidth=3.2cm, leftmargin=3.6cm]
  \item[Gleason tradition]
    Varadarajan, Kalmbach, Pt\'{a}k--Pulmannov\'{a}: $\sigma$-additive on
    a non-distributive OML, but real-valued on a fixed lattice --- not
    point-free.
  \item[D\"{o}ring 2009]
    \cite{Doering2009measures} point-free on the spectral presheaf over a
    directed poset --- but provably only \emph{finitely additive}, on a
    distributive Heyting algebra.
  \item[Bohrification (HLS)]
    \cite{HeunenLandsmanSpitters2009, HeunenLandsmanSpitters2011}
    point-free and the internal valuation \emph{is} $\sigma$-additive
    (Scott-continuous) --- but factors through an internally
    \emph{distributive} locale, $\sigma$-additive only per commutative
    context. Non-distributivity is tamed, not carried natively.
  \item[Localic valuations]
    Vickers, Simpson, Coquand--Spitters: point-free $\sigma$-additive, but
    on \emph{frames} --- distributive by definition.
  \item[OML Stone dualities]
    McDonald--Bimb\'{o}~\cite{McDonaldBimbo2023} and
    Cannon--D\"{o}ring~\cite{Cannon2013}: purely structural, no measure,
    no $\sigma$-version.
\end{description}
The unoccupied conjunction --- point-free $\times$ $\sigma$-additive
$\times$ natively non-distributive $\times$ directed-system-built --- is
exactly the descent residue. (The
directed-context-yields-non-distributivity \emph{idea} is itself not new;
it is the shape of Bohrification and of colimit-over-Boolean-contexts
generation of OMLs. Novelty would lie in the
$\sigma$-additive-point-free-native combination.)

\begin{remark}[Both standing negatives stress-tested]\label{rem:miss}
Two load-bearing negatives were re-checked against primary sources
(2026-06-04, both survived). \emph{(i)~No impossibility theorem} for the
\emph{concrete} class: the known no-go results are finite or two-valued,
and the one positive occupant $P(H)$$+$Gleason is \emph{non-concrete}
(Kochen--Specker~\cite{KochenSpecker1967}: no separating two-valued
states) --- so ``concrete'' is the load-bearing qualifier. \emph{(ii)~No
$\sigma$-Loomis--Sikorski}: both OML dualities take infinite joins as a
\emph{closure}, not a union (Cannon--D\"{o}ring $\mathrm{cls}(\bigcup
S_i)$; McDonald--Bimb\'{o} $(\bigcup h(a_n))^{\perp\perp}$). \emph{Guard:}
Gudder concrete logics \emph{are} set-representable non-Boolean --- but
point-ful posets, not a $\sigma$-LS theorem.
\end{remark}

% ==================================================================
\section{The open problem}
\label{sec:problem}
% ==================================================================

\begin{problem}
Does a $\sigma$-complete, concrete, non-Boolean, infinite OML admit a
Loomis--Sikorski-type representation (a $\sigma$-tribe of sets modulo a
$\sigma$-ideal), or a countable-join-preserving $\sigma$-Stone duality
--- or can one prove that none exists?
\end{problem}

\noindent A positive answer is the engine for \emph{relational probability
without realizations}: $\sigma$-additive probability built from the
entailment relation of a non-distributive OML, not descended from any
sample space --- the non-Boolean analogue of localic (pointless) measure
theory. A negative answer (an impossibility theorem) closes the whole
cluster, including both residues, at once. The outcome is binary and
large. No impossibility theorem is known for the live class
(Remark~\ref{rem:miss}); all three constructive routes are blocked
(\S\ref{sec:routes}). The contribution sharpens ``needs a new idea'' to
``needs a representation bypassing all three blocked routes.''

\begin{remark}[Why $L(H)$+Gleason does \emph{not} already settle this ---
the (A)/(B) point-space equivocation]\label{rem:AB}
The motivating idea (``probability from relations, point-free'') may
\emph{look} delivered by $L(H)$ + Gleason. It is not, and the reason is a
genuine equivocation between two point-spaces. \textbf{(A)} the Hilbert
rays of $H$, from which the Gleason density $\rho$ in $s = \tr(\rho\,\cdot)$
is built; \textbf{(B)} the McDonald--Bimb\'{o} dual filters $P(A)$. The
programme defines \emph{realization} as concentration on \textbf{(B)}.
Gleason removes \textbf{(B)} but \emph{requires} \textbf{(A)} --- and is a
\emph{representation} theorem, reducing every lattice state to the point
datum $\rho$, the most point-ful result available. So $L(H)$/Gleason
witnesses the $\PRlat / \PRdual$ \emph{separation} (a $\sigma$-additive
measure exists on the lattice while none concentrates on $P(A)$), not
point-freeness, and cannot serve as the existence proof sought. Because it
cannot, the $\sigma$-OML representation above is the \emph{only} candidate
engine. One must \emph{not} write ``Gleason already builds relational
probability.'' (Companion: \texttt{oml\_two\_point\_spaces.md}.)
\end{remark}

% ==================================================================
\section{The first concrete move}
\label{sec:first}
% ==================================================================

The computational and literature work is exhausted; what remains is human
mathematics, in two branches sharing one wall.
\begin{description}[style=nextline, font=\normalfont\bfseries,
    labelwidth=2.6cm, leftmargin=3cm]
  \item[(a)~The deep unlock]
    Build the Loomis--Sikorski-type / countable-join-preserving
    $\sigma$-Stone duality for a concrete non-Boolean infinite
    $\sigma$-OML, or prove none exists. High ceiling, no current foothold.
  \item[(b)~The tractable special case]
    The singular sliver: does a \emph{singular} orthoadditive state on
    $L(H)$ extend to a charge on $S_0(L(H))$? A concrete operator-algebra
    question (Calkin, Bunce--Wright, Takesaki), \emph{conditionally}
    self-contained.
\end{description}

\noindent\textbf{Attack (b) first} --- it is informative either way. If a
singular state fails to extend, then \emph{no} state on $L(H)$ extends, so
$\PRdual$ fails throughout $L(H)$; since the singular sliver is the only
candidate witness for ``descent does work independent of extension,'' this
closes that third-grade question negatively. If a singular state
\emph{does} extend, it is the first concrete object that extends but may
fail to concentrate, handing branch~(a) its motivating example.

\begin{problem}[The hinge that gates (b) --- settle this first]
Is the singular-extension obstruction reached with a \emph{finite} or a
\emph{countable} meet-zero family?
\end{problem}

\noindent This is \textbf{open}, and it decides whether (b) is genuinely
the easier branch. Proposition~\ref{prop:lh} is finitary (lines in a
$2$-plane) and so dodges the descent-axis wall. For singular states the
relevant pairwise-meet-zero families are \emph{infinite-dimensional}
subspaces with trivial intersection, so that \emph{line}-clustering route
is vacuous ($s(e)=0$ on all finite-dim $e$). But ``the line route is
vacuous'' is \emph{not} ``the obstruction is countable'': a \emph{finite}
family of infinite-dimensional pairwise-meet-zero subspaces with
$\sum_i s(a_i) > 1$ is unanalysed. If such a finite obstruction exists, (b)
is a self-contained operator-algebra problem to settle ahead of (a); if
only a countable family obstructs, (b) is entangled with the same
finitary-to-$\sigma$ bridge as (a). The very first thing to write down: a
finite family of infinite-dimensional, pairwise-trivially-intersecting
subspaces of $H$, and the question of whether orthoadditivity of a singular
state forces $\sum s > 1$. (Setup:
\texttt{verification/lh\_singular\_dichotomy.md}, the singular-state
constructions and the open statement.)

\begin{remark}[Methodological]
The thread's hard-won rules: verify by building the smallest example and
reading the primary source, not by recollection; weaken
``settled/strict/resolves'' to ``open/separation/conditional'' unless both
directions are checked; blocked routes $\neq$ impossible. A prior
``settled $=$ infinitary'' claim about the hinge above was an over-read of
the dichotomy file (which calls the question open) and was withdrawn --- a
cautionary instance of exactly the third rule.
\end{remark}

\begin{thebibliography}{99}
\small
\bibitem{BunceWright1992} L.~J.\ Bunce and J.~D.~M.\ Wright,
  \emph{The Mackey--Gleason problem}, Bull.\ Amer.\ Math.\ Soc.\ \textbf{26}
  (1992), 288--293.
\bibitem{BuresovaPtak2024} D.~Bure\v{s}ov\'{a} and P.~Pt\'{a}k,
  \emph{Set representation of orthomodular posets}, arXiv:2401.13798
  (2024).
\bibitem{Cannon2013} S.~Cannon and A.~D\"{o}ring,
  \emph{A generalisation of Stone duality to orthomodular lattices},
  in \emph{Reality and Measurement in Algebraic Quantum Theory},
  Springer (2018).
\bibitem{ChetcutiDvurecenskij2005} E.~Chetcuti and A.~Dvure\v{c}enskij,
  \emph{Measures on lattice effect algebras}, (2005).
\bibitem{deFinetti1974} B.~de~Finetti, \emph{Theory of Probability},
  Wiley (1974).
\bibitem{Doering2009measures} A.~D\"{o}ring,
  \emph{Quantum states and measures on the spectral presheaf},
  Adv.\ Sci.\ Lett.\ \textbf{2} (2009); arXiv:0809.4847.
\bibitem{Dvurecenskij2000} A.~Dvure\v{c}enskij,
  \emph{Loomis--Sikorski theorem for $\sigma$-complete MV-algebras and
  $\ell$-groups}, J.\ Austral.\ Math.\ Soc.\ \textbf{68} (2000), 261--277.
\bibitem{Dvurecenskij2005} A.~Dvure\v{c}enskij,
  \emph{Loomis--Sikorski representation of monotone $\sigma$-complete
  effect algebras with RDP}, (2005).
\bibitem{Fremlin3} D.~H.\ Fremlin, \emph{Measure Theory, Vol.~3:
  Measure Algebras}, Torres Fremlin (2002).
\bibitem{Freytes2020} H.~Freytes,
  \emph{An equational theory of $\sigma$-complete orthomodular lattices},
  Soft Computing \textbf{24} (2020).
\bibitem{Gleason1957} A.~M.\ Gleason,
  \emph{Measures on the closed subspaces of a Hilbert space},
  J.\ Math.\ Mech.\ \textbf{6} (1957), 885--893.
\bibitem{Harding1991} J.~Harding,
  \emph{Orthomodular lattices whose MacNeille completions are not
  orthomodular}, Order \textbf{8} (1991), 93--103.
\bibitem{HeunenLandsmanSpitters2009} C.~Heunen, N.~P.\ Landsman, B.~Spitters,
  \emph{A topos for algebraic quantum theory}, Comm.\ Math.\ Phys.\
  \textbf{291} (2009); arXiv:0709.4364.
\bibitem{HeunenLandsmanSpitters2011} C.~Heunen, N.~P.\ Landsman, B.~Spitters,
  \emph{Bohrification}, arXiv:0909.3468 (2011).
\bibitem{HornTarski1948} A.~Horn and A.~Tarski,
  \emph{Measures in Boolean algebras}, Trans.\ Amer.\ Math.\ Soc.\
  \textbf{64} (1948), 467--497.
\bibitem{Kalmbach1983} G.~Kalmbach, \emph{Orthomodular Lattices},
  Academic Press (1983).
\bibitem{KochenSpecker1967} S.~Kochen and E.~P.\ Specker,
  \emph{The problem of hidden variables in quantum mechanics},
  J.\ Math.\ Mech.\ \textbf{17} (1967), 59--87.
\bibitem{Kolmogorov} A.~N.\ Kolmogorov,
  \emph{Foundations of the Theory of Probability}, Chelsea (1956).
\bibitem{mahon2025} P.~S.\ Mahon, \emph{Structure from Observation
  (Paper~I)}, manuscript (2025).
\bibitem{mahonOMLnote} P.~S.\ Mahon, \emph{The extension problem for
  orthomodular observation algebras}, companion research note (2026).
\bibitem{McDonaldBimbo2023} J.~McDonald and K.~Bimb\'{o},
  \emph{A Stone-type duality for orthomodular lattices},
  Math.\ Logic Quarterly (2023); arXiv:2208.07430.
\bibitem{Navara1992} M.~Navara,
  \emph{Regularity does not imply $\sigma$-additivity}, (1992).
\bibitem{Pitowsky1989} I.~Pitowsky, \emph{Quantum Probability ---
  Quantum Logic}, Lecture Notes in Physics \textbf{321}, Springer (1989).
\bibitem{Ptak1987} P.~Pt\'{a}k, \emph{Exotic logics}, (1987).
\bibitem{PtakPulmannova1991} P.~Pt\'{a}k and S.~Pulmannov\'{a},
  \emph{Orthomodular Structures as Quantum Logics}, Kluwer (1991).
\bibitem{Takesaki1979} M.~Takesaki, \emph{Theory of Operator Algebras~I},
  Springer (1979).
\end{thebibliography}

\end{document}
